\documentclass[../../book-main.tex]{subfiles}
\graphicspath{{\subfix{../..}}}

\begin{document}

\chapter{Entropy, Diffusion, and Denoising}

\section{Proof of Entropy Increasing}

In this section appendix we provide a proof of \Cref{thm:diffusion_entropy_increases}. For convenience, we restate it as follows. 

\begin{theorem}[Theorem \Cref{thm:diffusion_entropy_increases}, Restated]
    Suppose that \(\vx\) is a random variable on \(\R^{d}\), \(\vg \sim \dNorm(0, I)\) independently of \(\vx\), \(T \in [0, \infty]\) is a time horizon, \(\sigma \colon [0, T] \to \R_{\geq 0}\) is an increasing differentiable function of \(t\). Let \((\vx_{t})_{t \in [0, T]}\) be a stochastic process defined as 
    \begin{equation}
        \vx_{t} := \vx + \sigma_{t}\vg, \qquad \forall t \in [0, T].
    \end{equation}
    Then, for any \(t > 0\):
    \begin{enumerate}
        \item The differential entropy \(h(\vx_{t})\) exists and is \(> -\infty\).
        \item The differential entropy is increasing: \(\odv*{h(\vx_{t})}{t} > 0\).
    \end{enumerate}
\end{theorem}
In this setting we additionally define the notations \(p\) for the law of \(\vx\), \(p_{t}\) for the law of \(\vx_{t}\), \(\phi\) for the law of \(\vg\) (i.e., a standard Gaussian random variable), and \(\phi_{t}\) for the law of \(\sigma_{t}\vg\) (i.e., a Gaussian random variable with mean \(0\) and covariance \(\sigma_{t}^{2}I\)). What is a ``law''? It is simply a function which maps sets to the probability that the random variable is inside the set, i.e.,
\begin{equation}
    \underbrace{p_{t}(A)}_{\text{law}} := \Pr[\vx_{t} \in A].
\end{equation}
By an abuse of notation, for each random variable, if it has a density then we will use the same notation for the law and the density, identifying them with each other as follows. 
\begin{equation}
    \underbrace{\phi(A)}_{\text{law}} = \int_{A}\underbrace{\phi(x)}_{\mathclap{\text{density}}}\odif{x} = \int_{A}\underbrace{\frac{1}{(2\pi)^{d/2}}\e^{-\frac{1}{2}\norm{x}_{2}^{2}}}_{\text{also density}}\odif{x}.
\end{equation}
With this notation and using the chain rule, it is easy to see that 
\begin{equation}
    \phi_{t}(x) = \frac{1}{\sigma_{t}^{d}}\phi\rp{\frac{x}{\sigma_{t}}}.
\end{equation}

Now we begin the proof. To prove \Cref{thm:diffusion_entropy_increases}, we will break it up into several fundamental lemmas involving the transformation of \(p_{t}\) over time \(t\). 
% \begin{itemize}
%     \item Some lemmas will just establish some technicality (i.e., for the sake of mathematical pedantry). These lemmas only exist to be used in the proof of other lemmas. We will annotate these with ``Technicality''. They may also invoke or deal with mathematical content which is a little too advanced for the rest of the book. Readers may safely skip these lemmas and their proofs on a first read-through.
%     \item Other lemmas will prove a central claim of \Cref{thm:diffusion_entropy_increases}. We aim to make these lemmas and their proofs as accessible as possible.
% \end{itemize}


\begin{lemma}
    The random variable \(\vx_{t}\) has a density given by 
    \begin{equation}
        p_{t}(u) = \Ex\phi_{t}(u - \vx),
    \end{equation}
    where (recall) \(\vx\) is the original random variable with law \(p\).
\end{lemma}
\begin{proof}
    \phantom{}

    \textbf{Step 1: Establishing that \(\vx_{t}\) has a density.} The rough idea is that since we add Gaussian noise to obtain \(\vx_{t}\), we know that \(\vx_{t}\) is supported everywhere and thus is a continuous random variable with density. The formal proof of this part draws on some higher-level mathematics to sort out definitions, so feel free to skip it unless curious.

    To show that \(\vx_{t}\) has a density, we show that if \(A\) is a Borel-measurable set with Borel measure \(0\), then \(p_{t}(A) = 0\). Indeed, since \(\vx\) and \(\vg\) are independent,
    \begin{equation}
        \vx_{t} = \vx + \sigma_{t}\vg \implies p_{t} = \phi_{t} * p
    \end{equation}
    where \(*\) denotes convolution (of probability densities/measures), it holds 
    \begin{align}
        p_{t}(A)
        &= \int_{\R^{d}}\indvar(u \in A)\odif{p_{t}(u)} \\
        &= \int_{\R^{d}}\indvar(u \in A)\odif{(\phi_{t} * p)(u)} \\ 
        &= \int_{\R^{d}}\int_{\R^{d}}\indvar(x + y \in A)\odif{\phi_{t}(x)}\odif{p(y)} \\
        &= \int_{\R^{d}}\int_{x \in (A - y)}\phi_{t}(x)\odif{x}\odif{p(y)} \\ 
        &= \int_{\R^{d}}\underbrace{\bs{\int_{x \in (A - y)}\phi_{t}(x)\odif{x}}}_{= 0}\odif{p(y)} \\ 
        &= 0 
    \end{align}
    where \(A - y = \{a - y \colon a \in A\}\) is the Minkowski translation, \textit{not} the set difference. Then the underbraced equality is true because the Borel measurable sets are closed under translation, i.e., \(A - y\) is a Borel set, and the Borel measure is translation-invariant, i.e., since \(A\) has Borel measure \(0\) then \(A - y\) has Borel measure \(0\), so integrating any function over \(A - y\) yields \(0\). So this proves that \(p_{t}\) has a density (which, remember, we will also call \(p_{t}\)).

    \textbf{Step 2: Finding the density of \(\vx_{t}\).}

    To show that \(\vx_{t}\) has the appropriate density,  we need to show that if \(A\) is any Borel-measurable set, then
    \begin{equation}
        p_{t}(A) = \int_{A}\Ex[\phi_{t}(u - \vx)]\odif{u}.
    \end{equation}
    Indeed, we can obtain this by just expanding the convolution in a very similar direction:
    \begin{align}
        p_{t}(A)
        &= \int_{\R^{d}}\int_{\R^{d}}\indvar(x + y \in A)\odif{\phi_{t}(x)}\odif{p(y)} \\
        &= \int_{\R^{d}}\int_{\R^{d}}\indvar(x + y \in A)\phi_{t}(x)\odif{x}\odif{p(y)} \\
        &= \int_{\R^{d}}\int_{\R^{d}}\indvar(u \in A)\phi_{t}(u - y)\odif{u}\odif{p(y)} \\
        &= \int_{\R^{d}}\int_{\R^{d}}\indvar(u \in A)\phi_{t}(u - y)\odif{p(y)}\odif{u} \\
        &= \int_{A}\bs{\int_{\R^{d}}\phi_{t}(u - y)\odif{p(y)}}\odif{u} \\
        &= \int_{A}\Ex[\phi_{t}(u - \vx)]\odif{u}
    \end{align}
    as desired. The only contentious equality here is the change in the order of integration, which is  justified by Tonelli's theorem and the fact that the integrand is the product of a normal density (measurable and positive everywhere) and an indicator function (measurable and non-negative everywhere) hence measurable and non-negative.
\end{proof}

The above fact is somewhat remarkable at first sight: even for a completely irregular random variable \(\vx\) (say, a Bernoulli random variable, which does not have a density), its Gaussian smoothing admits a density for every (arbitrarily small) \(t > 0\).

\begin{lemma}
    For \(t > 0\), the differential entropy \(h(\vx_{t})\) exists and is \(> -\infty\).
\end{lemma}
\begin{proof}
    We use a classic yet tedious analysis argument. Since \(\vx_{t}\) has a density, we can write 
    \begin{equation}
        h(\vx_{t}) = -\int_{\R^{d}}p_{t}(u)\log p_{t}(u) \odif{u}.
    \end{equation}
    Accordingly, let \(g \colon \R^{d} \to \R\) be defined as 
    \begin{equation}
        g(u) := -p_{t}(u)\log p_{t}(u) \implies h(\vx_{t}) = \int_{\R^{d}}g(u)\odif{u}.
    \end{equation}
    As usual to bound the value of an integral in analysis, define 
    \begin{equation}
        g_{+}(u) := \max(g(u), 0), \quad g_{-}(u) := \max(-g(u), 0) \quad \implies \quad g = g_{+} - g_{-}\quad \text{and} \quad g_{+}, g_{-} \geq 0.
    \end{equation}
    Then 
    \begin{equation}
        h(\vx_{t}) = \int_{\R^{d}}g_{+}(u)\odif{u} - \int_{\R^{d}}g_{-}(u)\odif{u},
    \end{equation}
    and both integrals are guaranteed to be non-negative since their integrands are. In order to prove that \(h(\vx_{t})\) exists, we need to show that \(\int_{\R^{d}}g_{+}(u)\odif{u} < \infty\) or \(\int_{\R^{d}}g_{-}(u)\odif{u} < \infty\). However, in order to show that \(h(\vx_{t}) > -\infty\), we just need to prove that \(\int_{\R^{d}}g_{-}(u)\odif{u} < \infty\). So we do the latter as it also proves the former.

    To bound the integral of \(g_{-}\) we need to understand the quantity \(g_{-}\). When is \(g\) negative?
    \begin{equation}
        g(u) \leq 0 \iff p_{t}(u)\log p_{t}(u) \geq 0 \iff \log p_{t}(u) \geq 0 \iff p_{t}(u) \geq 1.
    \end{equation}
    Thus, it holds that 
    \begin{equation}
        g_{-}(u) = \indvar(p_{t}(u) \geq 1)\cdot (-g(u)) = \indvar(p_{t}(u) \geq 1)p_{t}(u)\log p_{t}(u).
    \end{equation}
    In order to bound the integral of \(g_{-}(u)\), in this case we are lucky enough to bound the function \(g_{-}(u)\) itself. Namely, notice that \(\e^{-\norm{x}_{2}^{2}} \leq 1\), so \(\phi \leq \frac{1}{(2\pi)^{d/2}}\), so 
    \begin{equation}
        \phi_{t} \leq \frac{1}{(2\pi\sigma_{t}^{2})^{d/2}} := C_{t},
    \end{equation}
    which blows up as \(t \to 0\) but is finite for all finite \(t\). Therefore 
    \begin{equation}
        p_{t}(u) = \Ex \phi_{t}(u - \vx) \leq \Ex C_{t} = C_{t}.
    \end{equation}
    Now there are two cases.
    \begin{itemize}
        \item If \(C_{t} < 1\), then \(p_{t}(u) < 1\), so the indicator is never \(1\), hence \(g_{-} = 0\) identically and its integral is also \(0\).
        \item If \(C_{t} \geq 1\), then \(\log C_{t} \geq 0\), so since the logarithm is monotonically increasing,
        \begin{align}
            \int_{\R^{d}}g_{-}(u)\odif{u}
            &= \int_{\R^{d}}\indvar(p_{t}(u) \geq 1)p_{t}(u)\log p_{t}(u)\odif{u}  \\ 
            &= \Ex[\indvar(p_{t}(\vx_{t}) \geq 1)\log p_{t}(u)]  \\ 
            &\leq \Ex[\indvar(p_{t}(\vx_{t}) \geq 1) \log C_{t}] \\ 
            &= \Pr[p_{t}(\vx_{t}) \geq 1]\log C_{t}.
        \end{align}
    \end{itemize}
    Hence we have \(\int_{\R^{d}}g_{-}(u)\odif{u} < \infty\), so the differential entropy \(h(\vx_{t})\) exists and is \(> -\infty\).
\end{proof}

Now to prove that the differential entropy is increasing over time, we will use a characterization of the time-evolution of the probability density \(p_{t}\). Such characterizations are called \textit{Fokker-Planck equations}. Namely, we will show that the Fokker-Planck equation of \(\vx_{t}\) (hence \(p_{t}\)) follows a reparameterized \textit{heat equation}. The heat equation is a famous PDE which describes the diffusion of heat through space. This intuitively should make sense, and paints a mental picture: as the time \(t\) increases, the probability from the original (perhaps tightly concentrated) \(\vx\) disperses across all of \(\R^{d}\) like heat radiating from a source in a vacuum. But we need to prove this connection.

Central to the heat equation is a second-derivative term called the \textit{Laplacian} \(\Delta\), operating on a twice-differentiable function \(f \colon \R \times \R^{d} \to \R\) by 
\begin{equation}
    \Delta_{x} f_{t}(x) = \tr(\nabla_{x}^{2}f_{t}(x)) = \sum_{i = 1}^{d}\pdv[order=2]{f_{t}}{x_{i}}(x).
\end{equation}
Given this definition, the usual heat equation is the PDE 
\begin{equation}
    \pdv{f_{t}}{t}(x) = \Delta_{x} f_{t}(x).
\end{equation}
Theoretically, the heat equation describes \textit{time derivatives} of the densities \(p_{t}\) as functions of their instantaneous \textit{spatial derivatives} (in \(\vx\)), and it is an extremely useful, concise encapsulation of many regularity properties of the smoothed densities \(p_{t}\).\footnote{In fact, with a suitable (generalized) definition of the Laplacian operator, the heat equation actualy holds true in extreme generality---far beyond the case of additive Gaussian noise \cite{Bakry2016-pl}.}

Accounting for our parameter \(\sigma\), by playing around with both sides we observe the following heat-like Fokker-Planck equation.
\begin{lemma}
    The density \(p_{t}\) obeys the following heat-like PDE:
    \begin{equation}\label{eq:fokker_planck_ve}
        \pdv{p_{t}}{t}(u) = \sigma_{t}\sigma_{t}^{\prime}\cdot \Delta_{u} p_{t}(u).
    \end{equation}
    where \(\sigma_{t}^{\prime} = \odv{\sigma_{t}}{t}\) (as in \Cref{app:optimization}).
\end{lemma}
\begin{proof}
    Everything here is smooth and integrable, and so we can exchange derivatives and expectations (integrals) freely. Hence 
    \begin{equation}
        \pdv{p_{t}}{t}(u) = \pdv*{\Ex[\phi_{t}(u - \vx)]}{t} = \Ex\rs{\pdv*{\phi_{t}}{t}(u - \vx)}
    \end{equation}
    and 
    \begin{equation}
        \Delta_{u} p_{t}(u) = \Delta_{u} \Ex[\phi_{t}(u - \vx)] = \Ex[\Delta_{u}\phi_{t}(u - \vx)],
    \end{equation}
    the last equality following from the chain rule. It therefore suffices to show that
    \begin{equation}
        \pdv{\phi_{t}}{t}(z) = \sigma_{t}\sigma_{t}^{\prime}\cdot \Delta_{z} \phi_{t}(z),
    \end{equation}
    because then taking expectations on both sides will yield the desired equation \eqref{eq:fokker_planck_ve}.
    

    We will calculate the left side and right side and show that they are equal. We have 
    \begin{equation}
        \pdv{\phi_{t}}{t}(z) = \phi_{t}(z)\cdot \frac{\sigma_{t}^{\prime}}{\sigma_{t}^{3}}\cdot\bp{\norm{z}_{2}^{2} - d\sigma_{t}^{2}}
    \end{equation}
    by a direct calculation. On the other hand, 
    \begin{align}
        \pdv{\phi_{t}}{z_{i}}(z)
        &= -\phi_{t}(z)\cdot\frac{z_{i}}{\sigma_{t}^{2}} \\ 
        \pdv[order=2]{\phi_{t}}{z_{i}}(z)
        &= \phi_{t}(z)\cdot\frac{z_{i}^{2}}{\sigma_{t}^{4}} - \phi_{t}(z)\cdot\frac{1}{\sigma_{t}^{2}} = \phi_{t}(z)\cdot\frac{1}{\sigma_{t}^{4}}\cdot\bp{z_{i}^{2} - \sigma_{t}^{2}} \\ 
        \Delta_{z}\phi_{t}(z)
        &= \sum_{i = 1}^{d}\pdv[order=2]{\phi_{t}}{z_{i}}(z) = \phi_{t}(z)\cdot\frac{1}{\sigma_{t}^{4}}\cdot\bp{\norm{z}_{2}^{2} - d\sigma_{t}^{2}}.
    \end{align}
    So we have 
    \begin{equation}
        \pdv{\phi_{t}}{t}(z) = \phi_{t}(z)\cdot \frac{\sigma_{t}^{\prime}}{\sigma_{t}^{3}}\cdot\bp{\norm{z}_{2}^{2} - d\sigma_{t}^{2}} = \sigma_{t}\sigma_{t}^{\prime}\cdot \Delta_{z}\phi_{t}(z).
    \end{equation}
    The heat-like equation is satisfied for \(\phi_{t}\) and hence \(p_{t}\), as desired.
\end{proof}

Now that we have this, we have a very simple proof for the second contention of \Cref{thm:diffusion_entropy_increases}.
\begin{lemma}
    For \(t > 0\), the differential entropy increases:
    \begin{equation}
        \odv*{h(\vx_{t})}{t} > 0.
    \end{equation}
\end{lemma}
\begin{proof}
     Because everything is smooth and integrable, we can swap the derivative with the integral and obtain
    \begin{align}
        \odv*{h(\vx_{t})}{t}
        &= -\odv*{\int_{\R^{d}}p_{t}(u)\log p_{t}(u)\odif{u}}{t} \\
        &= -\int_{\R^{d}}\bs{\odv*{\bc{p_{t}(u)\log p_{t}(u)}}{t}}\odif{u} \\
        &= -\int_{\R^{d}} \odv{p_{t}}{t}(u)\bs{1 + \log p_{t}(u)}\odif{u} \\ 
        &= -\sigma_{t}\sigma_{t}^{\prime}\int_{\R^{d}}\Delta_{u}p_{t}(u)\bs{1 + \log p_{t}(u)}\odif{u} \\
        &= \sigma_{t}\sigma_{t}^{\prime}\int_{\R^{d}}\ip{\nabla_{u} \log p_{t}(u)}{\nabla_{u}p_{t}(u)}\odif{u} \\
        &= \sigma_{t}\sigma_{t}^{\prime}\int_{\R^{d}}\frac{\norm{\nabla_{u}p_{t}(u)}_{2}^{2}}{p_{t}(u)}\odif{u} \\ 
        &> 0.
    \end{align}
    Here the last inequality is because \(\sigma_{t} > 0\) and \(\sigma_{t}^{\prime} > 0\) for \(t > 0\), and the transfer from Laplacian to inner product is a result of Green's first identity, itself a consequence of integration by parts, and the identity \(\Delta f = \nabla \cdot (\nabla f)\) where \(\nabla \cdot\) is the divergence operator.
\end{proof}

This proves \Cref{thm:diffusion_entropy_increases}.

\section{Differential Entropy of Low-Dimensional Distributions}

In this short appendix we provide a proof of the following fact.

\begin{theorem}\label{thm:degenerate_entropy_negative_infinity}
    Let \(S\) be a Borel-measurable closed subset of \(\R^{d}\) with \(0\) volume.\footnote{Formally this means that \(S\) has Borel measure \(0\).} Let \(\vx\) be a random variable whose support is \(S\). Then \(h(\vx) = -\infty\).
\end{theorem}
\begin{proof}
    We will show this for the case \(S\) is compact and \(\vx\) is uniform on \(S\); this setting captures the essential ideas without being too technical. The basic idea is to take an \(\eps\)-thickening of \(S\), say \(S_{\eps}\) defined as 
    \begin{equation}
        S_{\eps} = \bigcup_{x \in S}B_{\eps}(x)
    \end{equation}
    where \(B_{\eps}(x)\) is the (\(\ell^{2}\)-norm) ball of radius \(\eps\) centered at \(x\), depicted in \Cref{fig:entropy_eps_thickening}.
    \begin{figure}
        \centering
        \begin{tikzpicture}
            \def\radius{0.2cm} % Define epsilon value
            \def\curve{(0,0) .. controls (1,1.5) and (3,-0.5) .. (4,1)} % Define the curve S

            % Draw many overlapping circles along the curve to represent the thickening S_eps
            % Use a decoration to place marks (circles) along the path
            \path[decoration={markings, mark=between positions 0 and 1 step 0.02 with {
                    % Fill a circle at each mark; low opacity to show overlap
                    \fill[red!30, opacity=0.5, draw=none] (0,0) circle (\radius);
                }}, postaction={decorate}] \curve;

            % Draw the original curve S on top, thicker and in blue
            \draw[blue, thick] \curve node[right, blue] {$S$};

            % Choose an example point x on the curve S
            \coordinate (p_on_curve) at (2, 0.5); % Approximate point on the curve
            % Draw the specific ball B_eps(x) for this x, slightly darker/more opaque
            \fill[red!50, opacity=0.7] (p_on_curve) circle (\radius);
            % Mark the point x
            \fill[black] (p_on_curve) circle (1pt) node[below left] {$x$};
            % Add an arrow indicating the radius epsilon for this specific ball

            % Label the thickened region S_eps, placing the label appropriately
             \node[red, below right] at (3.5, -0.2) {$S_{\eps} = \bigcup_{x \in S} B_{\eps}(x)$};

        \end{tikzpicture}
        \caption{Illustration of the \(\eps\)-thickening \(S_\eps\) of a curve \(S \subseteq \R^{2}\).}
        \label{fig:entropy_eps_thickening}
    \end{figure}
    We will work with random variables whose support is \(S_{\eps}\), which is fully-dimensional, and take the limit as \(\eps \to 0\).

    Thus let \(\vx \sim \operatorname{Uni}(S)\) such that it does not have a density\footnote{w.r.t.~the Borel measure on \(\R^{d}\).} and \(\vx_{\eps} \sim \operatorname{Uni}(S_{\eps})\). Since \(S_{\eps}\) has positive volume, \(\vx_{\eps}\) has a density \(p_{\eps}\) equal to
    \begin{equation}
        p_{\eps}(u) = \indvar(u \in S_{\eps}) \cdot \frac{1}{\mathrm{Vol}(S_{\eps})}.
    \end{equation}
    Since \(\vx\) does not have a density, the way to make sense of \(h(\vx)\) is through the equality 
    \begin{equation}
        h(\vx) = \lim_{\eps \searrow 0}h(\vx_{\eps});
    \end{equation}
    we will show that the latter limit evaluates to \(-\infty\).

    Indeed, using the convention that \(0 \log 0 = 0\), it holds
    \begin{align}
        h(\vx_{\eps}) 
        &= -\int_{\R^{d}}p_{\eps}(x)\log p_{\eps}(x)\odif{x} \\ 
        &= -\int_{S_{\eps}}\frac{1}{\operatorname{Vol}(S_{\eps})} \log\rp{\frac{1}{\operatorname{Vol}(S_{\eps})}}\odif{x} \\ 
        &= \frac{\log(\operatorname{Vol}(S_{\eps}))}{\operatorname{Vol}(S_{\eps})}\int_{S_{\eps}}\odif{x} \\ 
        &= \log(\operatorname{Vol}(S_{\eps})).
    \end{align}
    Since \(S\) is compact \(\operatorname{Vol}(S_{\eps})\) is finite and tends to \(0\) as \(\eps \searrow 0\).\footnote{In the case where \(S\) is not compact and \(\vx\) is not uniform on \(S\), we define a more refined version of \(\vx_{\eps}\) whose density at \(x\) changes at different rates in the directions orthogonal to and aligned with \(S\) at \(x\). This density results in a different calculation which does not involve the volume of \(S_{\eps}\) but rather something like the effective volume taken up by \(\vx_{\eps}\), which is finite and vanishing as \(\eps \searrow 0\).} Thus
    \begin{equation}
        h(\vx) = \lim_{\eps \searrow 0}h(\vx_{\eps}) = \lim_{\eps \searrow 0}\log(\operatorname{Vol}(S_{\eps})) = -\infty,
    \end{equation}
    as desired.
\end{proof}


\end{document}